% !TeX root = ../analysisII.tex

\section{Serie numeriche}
        Le serie numeriche sono somme di successioni da $B \subset \mathbb{N}$ a $\mathbb{R}$.
        Un esempio potrebbe essere il seguente: si prenda in considerazione la successione $\{a_{k}\}_{k \in \mathbb{N}}$ data da $a_{i} = \displaystyle \frac{1}{2^{i}}$, la corrispondente serie numerica è allora
        \[
            \sum_{i = 0}^{\infty} a_{i} = \sum_{i = 0}^{\infty} \frac{1}{2^{i}}
        \]

        \begin{definition}
            Sia $\{a_{k}\}_{k \in A}$, con $A \subset \mathbb{N}$ illimitato, una successione infinita in $\mathbb{R}$, si definisce $n$-esima somma parziale istanziata in $k_{0} \in A$, o ridotta $n$-esima, la seguente espressione
            \[
                S_{k_{0}, n} = \sum_{k = k_{0} \in \mathbb{N}}^{n} a_{k}
            \]
            Spesso, se $k_{0}$ è il valore di origine della successione ($k_{0} = \min(A)$), allora la ridotta $n$-esima si denota semplicemente $S_{n}$
        \end{definition}

        Uno degli oggetti di interesse dello studio delle serie numeriche è il valore del limite, data una successione reale infinita, della $n$-esima somma parziale al tendere di $n$ all'infinito.
        Individuiamo dunque quattro casi specifici:
        \begin{enumerate}[label=(\arabic*)]
            \item Se $\displaystyle \lim_{n \rightarrow \infty} S_{n} = s \in \mathbb{R}$, allora la serie si dice convergente a $s$, che si dice somma della serie, cioè:
            \[
                \sum_{k = 0}^{\infty} a_{k} = s
            \]
            \item Se $\displaystyle \lim_{n \rightarrow \infty} S_{n} = +\infty$, allora la serie si dice divergente positivamente, cioè:
            \[
                \sum_{k = 0}^{\infty} a_{k} = +\infty
            \]
            \item Se $\displaystyle \lim_{n \rightarrow \infty} S_{n} = -\infty$, allora la serie si dice divergente negativamente, cioè:
            \[
                \sum_{k = 0}^{\infty} a_{k} = -\infty
            \]
            \item Se $\displaystyle \nexists \lim_{n \rightarrow \infty} S_{n}$, allora la serie si dice indeterminata.
        \end{enumerate}
        Se il limite di una somma parziale istanziata in un generico $k_{0} \in \mathbb{N}$ è finito/infinito, allora tutti i limiti di tutte le somme parziali istanziate in un qualsiasi $k_{1} \in \mathbb{N}$ avranno limite finito/infinito.

        \begin{theorem}[Criterio di Cauchy]
            Sia $\sum_{k \in A} a_{k}$ con $A \subset \mathbb{N}$ illimitato una serie numerica. Se $\sum_{k \in A} a_{k} = a \in \mathbb{R}$, allora
            \[
                \lim_{n \rightarrow \infty} a_{k} = 0
            \]
        \end{theorem}
        \begin{proof}[Dimostrazione]
            Si consideri l'espressione $\mathcal{D}_{n} = S_{n + 1} - S_{n}$ e il suo limite al tendere di $n$ all'infinito, tenuto conto del fatto che $\lim_{n \rightarrow \infty} S_{n + 1} = \lim_{n \rightarrow \infty} S_{n}$.
            Quindi
            \[
                \lim_{n \rightarrow \infty} \mathcal{D}_{n} = \lim_{n \rightarrow \infty} (S_{n + 1} - S_{n}) = \lim_{n \rightarrow \infty} S_{n + 1} - \lim_{n \rightarrow \infty} S_{n} = a - a = 0 \implies \lim_{n \rightarrow \infty} \mathcal{D}_{n} = \lim_{n \rightarrow \infty} a_{n + 1} = \lim_{n \rightarrow \infty} a_{n} = 0
            \]
        \end{proof}

        \begin{definition}
            Una serie $\sum_{k \in A} a_{k}$ si dice a termini positivi se la successione sottostante $\{a_{k}\}_{k \in A}$ è a termini positivi.
        \end{definition}

        \begin{theorem}
            Una serie $\sum_{k \in A} a_{k}$ a termini positivi converge oppure diverge positivamente.
        \end{theorem}
        \begin{proof}[Dimostrazione]
            Si consideri la $n$-esima somma parziale $S_{n}$ come funzione di $n$. Questa è monotona crescente in quanto, se $n > m$, si ha che $S_{n} = S_{m} + S_{m + 1, n} > S_{m}$ perchè $S_{m + 1, n} > 0$ essendo somma di termini positivi.
            Allora
            \[
                \lim_{n \rightarrow \infty} S_{n} > 0 \implies \lim_{n \rightarrow \infty} S_{n} \in \mathbb{R}^{+} \lor \lim_{n \rightarrow \infty} S_{n} = +\infty
            \]

        \end{proof}

        \begin{definition}
            Una serie $\sum_{k = k_{0}}^{\infty} a_{k}$ si dice assolutamente convergente se $\sum_{k = k_{0}}^{\infty} |a_{k}|$ converge.
        \end{definition}

        \begin{proposition}
            Sia $\sum_{k = k_{0}}^{\infty} a_{k}$ una serie, allora
            \[
                \left| \sum_{k = k_{0}}^{\infty} a_{k} \right| \leq \sum_{k = k_{0}}^{\infty} |a_{k}|
            \]
        \end{proposition}

        \begin{definition}
            Una serie $\sum_{k = k_{0}}^{\infty} a_{k}$ si dice semplicemente convergente se converge ma non converge assolutamente.
        \end{definition}


        \subsection{Serie telescopiche}
            Una serie telescopica è una serie numerica tale per cui la successione sottostante è data nella forma $a_{k} = b_{k} - b_{k+1}$, con $\{b_{j}\}_{j \in \mathbb{N}}$ altra successione numerica.
            Così definite, le $n$-esime somme parziali delle serie telescopiche sono
            \[ S_{0} = a_{0} = b_{0} - b_{1} \]
            \[ S_{1} = a_{0} + a_{1} = b_{0} - b_{1} + b_{1} - b_{2} \]
            \[ S_{2} = a_{0} + a_{1} + a_{2} = b_{0} - b_{1} + b_{1} - b_{2} + b_{2} - b_{3} \]
            \[ \vdots \]
            \[ S_{n} = b_{0} - b_{n+1} \]

            Di conseguenza, il valore del limite di una serie telescopica è dato dal valore del limite della successione $\{b_{k}\}_{k \in \mathbb{N}}$:
            \begin{enumerate}[label=(\arabic*)]
                \item Se $\displaystyle \lim_{n \rightarrow \infty} b_{n} = b \in \mathbb{R}$, allora $\displaystyle \sum_{k = 0}^{\infty} a_{k} = b_{0} - b \in \mathbb{R}$
                \item Se $\displaystyle \lim_{n \rightarrow \infty} b_{n} = +\infty$, allora $\displaystyle \sum_{k = 0}^{\infty} a_{k} = +\infty$
                \item Se $\displaystyle \lim_{n \rightarrow \infty} b_{n} = -\infty$, allora $\displaystyle \sum_{k = 0}^{\infty} a_{k} = -\infty$
                \item Se $\displaystyle \nexists \lim_{n \rightarrow \infty} b_{n}$, allora $\displaystyle \nexists \sum_{k = 0}^{\infty} a_{k}$
            \end{enumerate}



        \subsection{Serie a termini alterni}
            Una serie a termini alterni è una serie numerica della forma
            \[
                \sum_{k = k_{0}}^{\infty} (-1)^{k} b_{k} \quad \text{con} \; b_{k} > 0 \;\; \forall k > k_{0}
            \]
            Inoltre, è possibile notare come
            \[
                \sum_{k = k_{0}}^{\infty} (-1)^{k + 1} b_{k} = \sum_{k = k_{0}}^{\infty} (-1)(-1)^{k} b_{k} = - \sum_{k = k_{0}}^{\infty} (-1)^{k} b_{k}
            \]



        \subsection{Operazioni tra serie numeriche}
            \subsubsection{Moltiplicazione serie-scalare}
                Data una generica serie numerica $\sum_{k = k_{0}}^{\infty} a_{k} = s \in \mathbb{R} \cup \{+\infty, -\infty\}$ e uno scalare $\lambda \in \mathbb{R} \setminus \{0\}$, si ha che
                \[
                    \lambda \sum_{k = k_{0}}^{\infty} a_{k} =
                    \begin{cases}
                        \lambda s \in \mathbb{R} &, \; \text{Se} \; s \in \mathbb{R} \\
                        +\infty                  &, \; \text{Se} \; s = +\infty \; \land \; \lambda > 0 \\
                        +\infty                  &, \; \text{Se} \; s = -\infty \; \land \; \lambda < 0 \\
                        -\infty                  &, \; \text{Se} \; s = -\infty \; \land \; \lambda > 0 \\
                        -\infty                  &, \; \text{Se} \; s = +\infty \; \land \; \lambda < 0 \\
                    \end{cases}
                \]

            \subsubsection{Somma tra serie numeriche}
                Date due generiche serie numeriche $\sum_{k = k_{0}}^{\infty} a_{k} = a \in \mathbb{R} \cup \{+\infty, -\infty\}$ e $\sum_{k = k_{0}}^{\infty} b_{k} = b \in \mathbb{R} \cup \{+\infty, -\infty\}$, si ha che
                \[
                    \sum_{k = k_{0}}^{\infty} a_{k} + \sum_{k = k_{0}}^{\infty} b_{k} = \sum_{k = k_{0}}^{\infty} (a_{k} + b_{k}) = a + b \iff a + b \; \text{non è forma indeterminata}
                \]




        \subsection{Criteri di convergenza}
            \subsubsection{Criterio del confronto maggiorante}
                Date due successioni $\{a_{k}\}_{k = k_{0}}^{\infty}$ e $\{b_{k}\}_{k = k_{0}}^{\infty}$ tali per cui
                \begin{enumerate}[label=(\arabic*)]
                    \item $a_{k} \geq 0, \; b_{k} \geq 0, \quad \forall k \geq k_{0}$
                    \item $a_{k} \leq b_{k}, \quad \forall k \geq k_{0}$
                \end{enumerate}
                si ha che
                \begin{itemize}
                    \item Se $\sum_{k = k_{0}}^{\infty} b_{k}$ converge, allora $\sum_{k = k_{0}}^{\infty} a_{k}$ converge.
                    \item Se $\sum_{k = k_{0}}^{\infty} a_{k}$ diverge positivamente, allora $\sum_{k = k_{0}}^{\infty} b_{k}$ diverge positivamente.
                \end{itemize}

            \subsubsection{Criterio del confronto asintotico}
                Date due successioni $\{a_{k}\}_{k = k_{0}}^{\infty}$ e $\{b_{k}\}_{k = k_{0}}^{\infty}$ tali per cui
                \begin{enumerate}[label=(\arabic*)]
                    \item $a_{k} \geq 0, \; b_{k} \geq 0, \quad \forall k \geq k_{0}$
                    \item $a_{k} \sim b_{k}, \quad k \rightarrow \infty$
                \end{enumerate}
                si ha che le serie $\sum_{k = k_{0}}^{\infty} a_{k}$ e $\sum_{k = k_{0}}^{\infty} b_{k}$ hanno stesso comportamento, ovvero che se una converge, anche l'altra converge; se una diverge, anche l'altra diverge.

            \subsubsection{Criterio della radice}
                Data una successione $\{a_{k}\}_{k = k_{0}}^{\infty}$ tale per cui
                \[
                    \lim_{k \rightarrow \infty} \sqrt[k]{a_{k}} = l \in \mathbb{R}^{+} \cup \{+\infty\}
                \]
                si ha che
                \begin{itemize}
                    \item Se $l < 1$, allora $\sum_{k = k_{0}}^{\infty} a_{k}$ converge.
                    \item Se $l > 1$, allora $\sum_{k = k_{0}}^{\infty} a_{k}$ diverge positivamente.
                \end{itemize}

            \subsubsection{Criterio del rapporto}
                Data una successione $\{a_{k}\}_{k = k_{0}}^{\infty}$ a termini positivi tale per cui
                \[
                    \lim_{k \rightarrow \infty} \frac{a_{k + 1}}{a_{k}} = l \in \mathbb{R}^{+} \cup \{+\infty\}
                \]
                si ha che
                \begin{itemize}
                    \item Se $l < 1$, allora $\sum_{k = k_{0}}^{\infty} a_{k}$ converge.
                    \item Se $l > 1$, allora $\sum_{k = k_{0}}^{\infty} a_{k}$ diverge positivamente.
                \end{itemize}

            \subsubsection{Criterio del confronto integrale}
                Sia $f: [k_{0}, +\infty) \rightarrow [0, +\infty)$ una funzione continua e decrescente e sia $\{a_{k}\}_{k = k_{0}}^{\infty}$ una successione tale per cui $a_{k} = f(k)$, si ha che
                \begin{enumerate}[label=(\arabic*)]
                    \item $\displaystyle \sum_{k = k_{0}}^{\infty} a_{k}$ ha lo stesso comportamento di $\displaystyle \int_{k_{0}}^{\infty} f(x) dx$
                    \item $\displaystyle \sum_{k = k_{0} + 1}^{\infty} a_{k} \leq \int_{k_{0}}^{\infty} f(x) dx \leq \sum_{k = k_{0}}^{\infty} a_{k}$
                    \item Se la serie converge ad un valore $s \in \mathbb{R}^{+}$, allora $\displaystyle S_{n} + \int_{k_{0} + 1}^{\infty} f(x) dx \leq s \leq S_{n} + \int_{k_{0}}^{\infty} f(x) dx$
                \end{enumerate}

            \subsubsection{Criterio di Leibnitz}
                Sia $\{b_{k}\}_{k = 0}^{\infty}$ una successione a termini positivi tale per cui
                \begin{itemize}
                    \item $\displaystyle \lim_{k \rightarrow \infty} b_{k} = 0$
                    \item $\{b_{k}\}$ è decrescente
                \end{itemize}
                Allora si ha che
                \begin{enumerate}[label=(\arabic*)]
                    \item $\displaystyle \sum_{k = 0}^{\infty} (-1)^{k} b_{k} = s \in \mathbb{R}$
                    \item $|s - S_{n}| \leq b_{n + 1}$
                    \item $S_{2n + 1} \leq s \leq S_{2n}, \; \forall n \in \mathbb{N}$
                \end{enumerate}




    \section{Serie di potenze}
        Dato un $x_{0} \in \mathbb{R}$ e una successione reale $\{a_{k}\}_{0}^{\infty}$ si dice serie di potenze centrata in $x_{0}$ di coefficienti $\{a_{k}\}$ la seguente espressione
        \[
            \sum_{k = 0}^{\infty} a_{k}(x - x_{0})^{k} = a_{0} + a_{1}(x - x_{0}) + a_{2}(x - x_{0})^{2} + ...  \quad \text{con} \; x \in \mathbb{R}
        \]
        Si dice invece insieme di convergenza della serie l'insieme
        \[
            A = \{ x \in \mathbb{R} : \sum_{k = 0}^{\infty} a_{k}(x - x_{0})^{k} \; \text{converge} \; \}
        \]

        Si può osservare come, con un opportuno cambio di variabile, è possibile rendere tutte le serie centrate in un generico $x_{0} \in \mathbb{R}$ in serie centrate in $0$
        \[
            \sum_{k = 0}^{\infty} a_{k}(x - x_{0})^{k} = \sum_{k = 0}^{\infty} a_{k}t^{k} \quad \text{con} \; t = (x - x_{0})
        \]

        \begin{theorem}
            Sia $\sum_{k = 0}^{\infty} a_{k}x^{k}$ una serie di potenze centrata in $0$.
            \begin{enumerate}
                \item Se la serie converge per un $x_{1} \in \mathbb{R}$ diverso da $0$, allora la serie converge assolutamente per ogni $x \in \mathbb{R}$ tale per cui $|x| < |x_{1}|$
                \item Se la serie diverge per un $x_{2} \in \mathbb{R}$, allora la serie diverge per ogni $x \in \mathbb{R}$ tale per cui $|x| > |x_{2}|$
            \end{enumerate}
        \end{theorem}

        \begin{proof}[Dimostrazione]
            Se la serie converge in $x_1\neq0$, la successione dei suoi
            termini $a_kx_1^k$ tende a zero ed è quindi limitata: esiste
            $M>0$ tale che $\lvert a_kx_1^k\rvert\leq M$ per ogni $k$.
            Per $\lvert x\rvert<\lvert x_1\rvert$, posto
            $q=\lvert x\rvert/\lvert x_1\rvert<1$, si ha
            \[
                \lvert a_kx^k\rvert
                =\lvert a_kx_1^k\rvert q^k
                \leq Mq^k.
            \]
            La convergenza assoluta segue per confronto con la serie
            geometrica $\sum_{k=0}^{+\infty}Mq^k$.

            Per la seconda affermazione, supponiamo per assurdo che la
            serie converga in un punto $x$ con
            $\lvert x\rvert>\lvert x_2\rvert$. Per quanto appena dimostrato
            essa convergerebbe allora assolutamente in $x_2$, in
            contraddizione con l'ipotesi. Deve quindi divergere per ogni
            $x$ tale che $\lvert x\rvert>\lvert x_2\rvert$.
        \end{proof}

        \begin{definition}
            Data una serie di potenze $\sum_{k = 0}^{\infty} a_{k}(x - x_{0})^{k}$ si dice raggio di convergenza della serie il seguente
            \[
                R = \sup\{ x \in \mathbb{R} : \sum_{k = 0}^{\infty} a_{k}(x - x_{0})^{k} \; \textnormal{converge} \; \}
            \]
        \end{definition}

        A seconda del raggio di convergenza di una serie di potenze, si possono trarre le seguenti conclusioni
        \begin{enumerate}[label=(\arabic*)]
            \item $R = 0$, allora $A = \{x_{0}\}$.
            \item $R \in [0, +\infty)$, allora $A = (-R + x_{0}, R + x_{0})$. Gli estremi $\{-R + x_{0}\}$ ed $\{R + x_{0}\}$ possono essere inclusi solo dopo un'analisi approfondita circa la serie.
            \item $R = +\infty$, allora $A = \mathbb{R}$
        \end{enumerate}



        \subsection{Criteri di calcolo del raggio di convergenza}
            \subsubsection{Criterio della radice}
                Sia $\sum_{k = 0}^{\infty} a_{k}(x - x_{0})^{k}$ una serie di potenze.
                Se
                \[
                    \lim_{k \rightarrow \infty} \sqrt[k]{a_{k}} = l \in [0, +\infty]
                \]
                Allora
                \[
                    R = \begin{cases}
                        0, & l = +\infty \\
                        \displaystyle \frac{1}{l}, & l \in (0, +\infty) \\
                        +\infty, & l = 0
                    \end{cases}
                \]
                dove $R$ è il raggio di convergenza della serie.

            \subsubsection{Criterio della radice}
                Sia $\sum_{k = 0}^{\infty} a_{k}(x - x_{0})^{k}$ una serie di potenze tale per cui $a_{k} \neq 0, \; \forall k$.
                Se
                \[
                    \lim_{k \rightarrow \infty} \frac{|a_{k+1}|}{|a_{k}|} = l \in [0, +\infty]
                \]
                Allora
                \[
                    R = \begin{cases}
                        0, & l = +\infty \\
                        \displaystyle \frac{1}{l}, & l \in (0, +\infty) \\
                        +\infty, & l = 0
                    \end{cases}
                \]
                dove $R$ è il raggio di convergenza della serie.



        \subsection{Derivate e integrali di serie di potenze}
            \begin{definition}
                Sia $\sum_{k = 0}^{\infty} a_{k}(x - x_{0})^{k}$ una serie di potenze con raggio $R > 0$.
                Si dice serie derivata la serie
                \[
                    \left( \sum_{k = 0}^{\infty} a_{k}(x - x_{0})^{k}\right)' = \sum_{k = 1}^{\infty} a_{k}k(x - x_{0})^{k - 1}
                \]
            \end{definition}

            \begin{definition}
                Sia $\sum_{k = 0}^{\infty} a_{k}(x - x_{0})^{k}$ una serie di potenze con raggio $R > 0$.
                Si dice serie integrale la serie
                \[
                    \int_{x_{0}}^{x} \left( \sum_{k = 0}^{\infty} a_{k}(t - x_{0})^{k} \right) dt = \sum_{k = 0}^{\infty} a_{k}\frac{(x - x_{0})^{k + 1}}{k + 1}
                \]
            \end{definition}

            \begin{theorem}
                Sia $\sum_{k = 0}^{\infty} a_{k}(x - x_{0})^{k} = S(x)$ una serie di potenze con raggio $R > 0$.
                Allora
                \begin{enumerate}[label=(\arabic*)]
                    \item La serie $S(x)$ è di classe $C^{\infty}(x_{0} - R, x_{0} + R)$.
                    \item L'espressione della serie derivata di ordine $n \in \mathbb{N}$ è
                    \[
                        S^{(n)}(x) = \sum_{k = 0}^{\infty} \frac{d^{n}}{dx^{n}} \left[ a_{k}(x - x_{0})^{k} \right] = \sum_{k = n}^{\infty} a_{k} \frac{k!}{(k-n)!}(x - x_{0})^{k - n} , \; \forall x \in (x_{0} - R, x_{0} + R)
                    \]
                    \item $\displaystyle a_{n} = \frac{S^{(n)}(x_{0})}{n!}$
                    \item La serie integrale $\int S(x)$ ha lo stesso raggio di convergenza di $S(x)$
                    \[
                        \sum_{k = 0}^{\infty} a_{k}\frac{(x - x_{0})^{k + 1}}{k + 1} = \int_{x_{0}}^{x} S(t) dt, \; \forall x \in (x_{0} - R, x_{0} + R)
                    \]
                \end{enumerate}
            \end{theorem}
            Anche per le derivate e gli integrali di serie il comportamento circa gli estremi dell'insieme di convergenza va studiato caso per caso.



        \subsection{Funzioni analitiche}
        Una funzione $f \in C^{\infty}(x_{0} - \epsilon, x_{0} + \epsilon)$ con $\epsilon > 0$ si dice analitica in $x_{0}$ se
        \begin{enumerate}[label=(\arabic*)]
            \item $\displaystyle \sum_{k = 0}^{\infty} \frac{f^{(k)}(x_{0})}{k!}(x - x_{0})^{k}$ ha raggio di convergenza $R$ non nullo
            \item $\displaystyle \sum_{k = 0}^{\infty} \frac{f^{(k)}(x_{0})}{k!}(x - x_{0})^{k} = f(x)$ per ogni $x \in (x_{0} - \delta, x_{0} + \delta), \; R > \delta > 0$
        \end{enumerate}

        L'espressione $\displaystyle \sum_{k = 0}^{\infty} \frac{f^{(k)}(x_{0})}{k!}(x - x_{0})^{k}$ è chiamata serie di Taylor (qualora $x_{0} = 0$ essa è anche detta serie di MacLaurin) di $f$ centrata in $x_{0}$

        \begin{theorem}[Condizioni sufficienti per l'analiticità]
            Sia $f \in C^{\infty}(x_{0} - \epsilon, x_{0} + \epsilon)$ con $\epsilon > 0, \; x_{0} \in \mathbb{R}$.
            Se $\exists M > 0$ tale per cui $f^{(k)}(x) \leq M$ per ogni $x \in (x_{0} - \epsilon, x_{0} + \epsilon), \; k \in \mathbb{N}$, allora $f$ è analitica in $x_{0}$ e il raggio di convergenza della serie di Taylor di $f$ centrata in $x_{0}$ è $\geq \epsilon$.
        \end{theorem}



    \section{Serie di Fourier}
    Sia $\mathcal{P} \subset L^{2}([0, T])$ lo spazio vettoriale di tutte le funzioni periodiche di periodo $T$ e continue a tratti su $[0, T]$.
    E' possibile dimostrare che
    \[
        \mathcal{B} = \left\{ \left. 1, \cos(n\omega_{0}t), \sin(n \omega_{0} t) \;  \right| \; n \geq 1, \; \omega_{0} = \frac{2\pi}{T} \right\}
    \]
    è una base ortogonale dello spazio $\mathcal{P}$.
    Data una funzione $f \in \mathcal{P}$, è possibile esprimerla secondo la base $\mathcal{B}$
    \[
        f(t) = a_{0} + \sum_{n = 1}^{\infty} \left[ a_{n}\cos(n \omega_{0} t) + b_{n}\sin(n \omega_{0} t) \right]
    \]
    dove i coefficienti sono dati da
    \[
        a_{0} = \frac{\langle f, 1\rangle}{||1||^{2}} = \frac{1}{T} \langle f, 1 \rangle = \frac{1}{T} \int_{0}^{T} f(t) \cdot 1 \: dt = \frac{1}{T} \int_{0}^{T} f(t) dt
    \]
    \[
        a_{n} = \frac{\langle f, \cos(n \omega_{0} t) \rangle}{||\cos(n \omega_{0} t)||^{2}} = \frac{2}{T} \langle f, \cos(n \omega_{0} t) \rangle = \frac{2}{T} \int_{0}^{T} f(t)\cos(n \omega_{0} t) dt
    \]
    \[
        b_{n} = \frac{\langle f, \sin(n \omega_{0} t) \rangle}{||\sin(n \omega_{0} t)||^{2}} = \frac{2}{T} \langle f, \sin(n \omega_{0} t) \rangle = \frac{2}{T} \int_{0}^{T} f(t)\sin(n \omega_{0} t) dt
    \]
    con $\langle f, g \rangle = \int_{0}^{T}fg$ prodotto scalare in $\mathcal{P}$.
    L'espressione
    \[
        f(t) = a_{0} + \sum_{n = 1}^{\infty} \left[ a_{n}\cos(n \omega_{0} t) + b_{n}\sin(n \omega_{0} t) \right]
    \]
    è detta serie di Fourier di $f$.

    \begin{proposition}
        Sia $f \in \mathcal{P}$, allora
        \begin{itemize}
            \item Se $f$ è pari, allora $b_{n} = 0, \; \forall n \geq 1$
            \item Se $f$ è dispari, allora $a_{n} = 0, \; \forall n \geq 1$
        \end{itemize}
    \end{proposition}

    \begin{proof}[Dimostrazione]
        Si consideri il caso in cui $f$ sia pari.
        Si ha che $f \sin(n\omega_{0}t)$ è una funzione dispari.
        In aggiunta, vale la seguente relazione per periodicità di $f$ e $\sin(n\omega_{0}t)$
        \[
            \frac{2}{T} \int_{0}^{T} f(t)\sin(n \omega_{0} t) dt = \frac{2}{T} \int_{-\frac{T}{2}}^{\frac{T}{2}} f(t)\sin(n \omega_{0} t) dt
        \]

        Mettendo il tutto insieme, si ha
        \[
            b_{n} = \frac{2}{T} \int_{0}^{T} f(t)\sin(n \omega_{0} t) dt = \frac{2}{T} \int_{-\frac{T}{2}}^{\frac{T}{2}} f(t)\sin(n \omega_{0} t) dt = 0
        \]

        Nel caso in cui $f$ fosse dispari, allora $f \cos(n \omega_{0} t)$ sarebbe dispari, quindi
        \[
            a_{n} = \frac{2}{T} \int_{0}^{T} f(t)\cos(n \omega_{0} t) dt = \frac{2}{T} \int_{-\frac{T}{2}}^{\frac{T}{2}} f(t)\cos(n \omega_{0} t) dt = 0
        \]
    \end{proof}

    \begin{theorem}[Criterio di convergenza della serie di Fourier]
        Sia $f \in \mathcal{P}$ e sia $A = \{t_{i}\}_{i \in \mathbb{N}}$ l'insieme delle discontinuità di $f$ in $[0, T]$.
        Se esistono i limiti $\displaystyle \lim_{t \rightarrow t_{i}} f(t)$ e $\displaystyle \lim_{t \rightarrow t_{i}} f'(t)$ e $f \in C^{1}(t_{i}, t_{i + 1})$ per ogni $t_{i} \in A$, allora per ogni $t_{0} \in \mathbb{R}$ la serie di Fourier converge a
        \[
            \frac{1}{2} \left[ \lim_{t \rightarrow t_{0}^{+}} f(t) + \lim_{t \rightarrow t_{0}^{-}} f(t) \right]
        \]
    \end{theorem}

    \begin{theorem}[Teorema della convergenza dell'energia]
        Sia $f \in \mathcal{P}$. Allora
        \[
            ||f||^{2}_{2} = \int_{0}^{T} f(t)^{2} dt = a_{0}^{2} \: T + \frac{T}{2} \: \sum_{n = 1}^{\infty} \left[a_{n}^{2} + b_{n}^{n}\right]
        \]
    \end{theorem}

    \subsection{Serie di fourier di funzioni complesse}
    Sia $\mathcal{C} \subset L^{2}([0, T], \mathbb{C})$ lo spazio vettoriale di tutte le funzioni complesse a variabile reale con periodo $T$.
    E' possibile dimostrare che
    \[
        \mathcal{B} = \left\{ e^{ik\omega_{0}t} \; | \; k \in \mathbb{Z} \right\}
    \]
    è una base di $\mathcal{C}$.
    Di conseguenza, data una funzione $f \in \mathcal{C}$, si ha
    \[
        f(t) = \sum_{k \in \mathbb{Z}} c_{k} e^{ik\omega_{0}t}
    \]
    dove
    \[
        c_{k} = \frac{\langle f, e^{ik\omega_{0}t} \rangle}{||e^{ik\omega_{0}t}||^{2}} = \frac{1}{T} \langle f, e^{ik\omega_{0}t} \rangle  = \frac{1}{T} \int_{0}^{T} f(t) \overline{e^{ik\omega_{0}t}} dt
    \]
    \[
        ||f||^{2}_{2} = \int_{0}^{T} f(t) \overline{f(t)} dt = T \sum_{k \in \mathbb{Z}} |c_{k}|^{2}
    \]
